\par\noindent\begin{minipage}{\linewidth}
% !TeX root = ../教材题库_全部条目_转写工作稿.tex
% Source: XJ-B1-C02-CA-P033-Y0234
% 第2章 / 2.1 相等关系与不等关系 / 教材 p.33 / PDF 第40页
% 已对照原页转写题干及公式；待教师复核；未制作答案
% 原题及相邻原解截图（核对用）：../题目原图/XJ-B1-C02-CA-P033-Y0234-01.webp
% 本文件只在转写完成后作为题干使用；OCR 条目尚不是可靠题干。
\begin{bankitem}{2.1 案例 案例1}{第2章 / 2.1 相等关系与不等关系 / 教材 p.33 / PDF 第40页}{XJ-B1-C02-CA-P033-Y0234}{已对照原页转写题干及公式；待教师复核；未制作答案}
2002 年在北京召开的第 24 届国际数学家大会的会标如图 2.1-2（1）所示，其中的几何图形取材于我国古代数学家对勾股定理的证明，如图 2.1-2（2）。图中的正方形 $ABCD$ 由 4 个全等的直角三角形和小正方形 $EFGH$ 组成。设直角三角形的两条直角边长分别为 $a,b\ (a\ne b)$，那么正方形 $ABCD$ 的边长为 $\sqrt{a^2+b^2}$。
\bankfig{C02-CA033}{0.62}
又 4 个直角三角形的面积之和为 $4\times\frac12ab=2ab$，正方形 $ABCD$ 的面积为 $a^2+b^2$，并且正方形 $ABCD$ 的面积大于 4 个直角三角形的面积之和，于是可得不等式 $a^2+b^2>2ab\ (a\ne b)$。
\end{bankitem}
\end{minipage}\par\medskip
